\newpage
\section{Related Work}

Chapter 2 established what the evidence indicates should be effective. This chapter examines what has been built. Section 3.1 reviews research prototypes; Section 3.2 reviews deployed systems, internationally and in Vietnam; Section 3.3 derives comparison criteria from the attributes the background identifies as consequential for learning systems, applies them, and states the resulting gap. Throughout, attention is given both to the limitations of existing work and to the design solutions worth adopting from it.

\subsection{Related work from research literature}

A growing body of research applies games and generative models to career development. It establishes the direction as active and credible, and three limitations recur across it.

Reflection as the terminal outcome. The most developed prototypes assist users in thinking about a career rather than performing one. CareerSim\cite{r76} is a language-model-driven career development game combining generated content with alumni data and career construction theory, advancing an avatar through life stages to produce a personalised report; a preliminary study with twelve participants reported improved understanding and reflection. Its interaction concerns a career trajectory rather than the tasks constituting an occupation, and the authors note limited contextual depth and difficulty with distant scenarios. Future You\cite{r77}, in which participants conversed with a model-generated future self, reduced anxiety and increased future self-continuity across 344 participants, and Letters from Future Self\cite{r78} found comparable benefits for career exploration and resilience with thirty-six participants. Both demonstrate measurable psychological benefit from structured imagination of a future self. Neither provides the basis on which a user might discover, through performance, that the imagined self would find the work unsuitable --- the informational deficit identified in Section 2.1. The design element worth adopting is the prospective framing itself, which these studies show is motivationally effective and which can be applied at the entry point of an experiential system rather than as its substance.

Trajectory abstraction without task fidelity. A second group models career development under uncertainty while abstracting from the work. CareerPooler\cite{r79} represents career development through a physical metaphor in which generated events perturb a trajectory, and reports significant gains in engagement, information gain, satisfaction and career clarity relative to a conversational baseline across twenty-four participants. It constitutes the strongest available evidence that generated, interactive career exploration outperforms conversational advice, and its treatment of unplanned events accords with the theoretical position in Section 2.4. Because the representation is metaphorical, however, the user does not perform occupational tasks, and outcomes are consequently reported in terms of clarity rather than demonstrated capability. Its event-generation pipeline, which balances favourable and unfavourable outcomes to avoid a uniformly positive trajectory, is directly adaptable to any system incorporating stochastic career events. Comparable abstraction characterises generative career imagination for children\cite{r80} and career simulation for self-regulation in secondary students\cite{r81}.

Guidance without measurement. A third group improves the guidance interaction while leaving assessment unaddressed. A career counselling agent grounded in self-determination theory\cite{r82} reduced decision-making difficulty and improved engagement, confirming the applicability of the motivational framework in Section 2.3; its outcomes are measured through self-report instruments rather than through anything the user demonstrably did. Among the systems reviewed here, none infers competence from behaviour within a task, and none validates automated judgment against an expert-constructed standard. The techniques required are established outside the career domain --- memory-consistent interactive narrative\cite{r53,r56}, rubric-based automated evaluation with its attendant validation requirements\cite{r57,r59,r60}, and unobtrusive behavioural assessment\cite{r75} --- but have not been combined within it.

Evaluation practice in this literature is a further limitation. Sample sizes range from twelve to thirty-six, with one exception at 344, and designs are predominantly qualitative or within-subject. Such designs are appropriate for establishing feasibility and eliciting user response, but they do not support claims about assessment accuracy, which requires comparison against an external standard.

\subsection{Related work from deployed systems}

\subsubsection{International systems}

Authored task simulation. Forage\cite{r83} is the established platform in experiential career exposure: employer-authored simulations in which users attempt representative tasks and compare their work against a model answer, distributed without charge and funded through employer talent acquisition. It substantiates the underlying premise that users will voluntarily attempt occupational tasks, and its task decomposition --- discrete, self-contained work products with supporting context --- is a directly adaptable model for scenario structure. Its constraints follow from its authoring model. Content expands only as employers commission it; repetition yields identical material, limiting exploration to a single pass; and assessment terminates in a fixed model answer, which demonstrates what a competent response looks like without evaluating the response the user produced. The catalogue also reflects its funding, concentrating in fields that sponsor recruitment. Human-mediated alternatives such as Virtual Internships and Extern provide genuine workplace exposure at correspondingly limited scale, while career video and mentoring platforms provide exposure without performance.

Generated task simulation and workplace assessment. A more recent group generates content rather than commissioning it. EntryLevel\cite{r84} produces AI-generated job simulations parameterised by field, skill and difficulty with an accompanying automated tutor, establishing commercially that occupational task content can be generated; it delivers these as discrete projects without progression or a connecting structure between occupations. Anthropos\cite{r85} conducts extended workplace scenarios in which a user interacts with several model-driven colleagues and produces work artefacts, explicitly generating behavioural evidence rather than self-report ratings. It is the closest deployed analogue to behavioural assessment as described in Section 2.2, and its multi-role scenario construction is adaptable. Its orientation is nonetheless towards personnel selection: the evidence generated serves an employer's hiring decision rather than a user's exploration, which determines both what is measured and who receives it. A third group --- including conversation and interview practice systems with adaptive personas and scored feedback\cite{r86} --- rehearses the selection process rather than the work, and its scoring dimensions, being conversational, do not generalise to occupational task performance.

Generative narrative systems. The systems that have most successfully solved the generation problems described in Section 2.5 are not career systems. AI Dungeon\cite{r87} and character.ai\cite{r88} established substantial demand for model-driven interactive narrative, particularly among the demographic most relevant to career exploration, and equally demonstrate the consistency degradation that unconstrained generation produces at scale. Hidden Door\cite{r89} is the more instructive case: it generates narrative within bounded authored components rather than through open generation, with the reported consequence that its narratives remain coherent across sessions. This is independent commercial corroboration of the finding in Section 2.5 that authored structure, not model capability, governs consistency, and its architecture is the most directly adaptable element identified in this review. What these systems lack is any representational commitment to a real domain, any model of competence, and any assessment.

Assessment and orientation instruments. O*NET and its associated public interfaces, together with commercial and institutional counterparts, constitute the most widely used category. They administer an inventory, return ranked occupational suggestions, and terminate. Their occupational databases are a genuine asset --- and, as noted in Section 2.6, a usable one --- but the guidance model rests on the self-report foundation whose limitations Section 2.4 describes.

\subsubsection{Vietnamese systems}

The domestic market is materially homogeneous, which is itself the relevant finding.

Assessment-based provision. JobWay\cite{r90}, reported in Vietnamese government press as the first digitised career orientation application in the country, comprises interest and personality assessment, descriptions of approximately three hundred occupations drawn from international sources, and linked institutional information. Notably, its occupational descriptions include the common difficulties of each occupation --- an implicit application of the realistic job preview principle, delivered as text. JobTest.vn\cite{r91} is the most instrumentally developed, offering multiple validated inventories alongside an aptitude instrument adapted and standardised for Vietnamese use and an occupational database of substantial scale; its sophistication is concentrated in the measurement of self-report. Hướng nghiệp Sông An\cite{r92} combines counsellor training with a library of licensed instruments, representing the most methodologically careful domestic provision and remaining entirely assessment- and counsellor-based. Hướng Nghiệp LwL\cite{r93} advertises experiential activity alongside assessment and occupational information for a younger audience, though its substance remains assessment, content and mentoring.

Generative systems in the domestic market. Recent entrants apply generative models to career guidance, uniformly in conversational form. AI Hay\cite{r94} provides a model-driven agent predicting admission likelihood and recommending institutions and fields, together with an automated career-personality instrument; a parallel development across Vietnamese universities has produced admissions advisory agents at institutional scale. These establish market acceptance of AI-assisted career guidance domestically. None provides performed experience of occupational work, and none assesses demonstrated capability. Among the systems surveyed, no Vietnamese system offers generated, performed occupational simulation --- a claim that should be re-verified close to submission, given the rate of change in this market.

\subsection{Comparison, gap analysis, and positioning}

\subsubsection{Derivation of comparison criteria}

Comparing systems against the features of a proposed system produces a circular result. The criteria below are instead derived from attributes that Chapter 2 identifies as determining whether a system of this general class --- one intended to help a person learn about and choose an occupation --- achieves its purpose. Each is stated as a general property, applicable to any system in the category, and each traces to an evidential basis rather than to a design preference.

The first concerns the nature of what the user does. Experiential learning and simulation transfer research (2.1) indicate that engagement with representative tasks of actual practice conveys understanding that description does not; task authenticity is therefore a primary attribute. The second concerns what is conveyed. The realistic job preview literature (2.1) and the decision-difficulty taxonomy (1.1) establish that balanced information, including unattractive aspects, is what improves fit --- hence completeness and candour of occupational representation. The third concerns responsiveness: adaptive instruction and the generative literature (2.5) indicate that content which responds to the individual outperforms fixed content, giving adaptivity to the learner. The fourth concerns whether the system knows what the user can do: evidence-centered design (2.2) establishes both the possibility of inferring competence from behaviour and the conditions under which such inference is valid, giving validity of competence measurement. The fifth follows from it --- assessment is educationally inert without actionable response, giving quality of feedback and guidance. The sixth derives from the gamification and self-determination literatures (2.3): sustained engagement determines whether any learning occurs at all. The seventh concerns navigation: career theory and competency-based progression (2.2, 2.4) indicate that occupational choice involves comparing and moving among related options, giving support for pathway navigation. The eighth concerns durability: the authoring literature (2.5) shows that content about real occupations degrades unless it can be extended and maintained, giving scalability and currency of content. A ninth attribute, the existence of empirical validation, is treated in the discussion rather than tabulated, since it concerns the evidence for a system rather than a property of it.

\subsubsection{Comparative assessment}

Ratings use three levels, defined as follows. Full ($\bullet$) indicates the attribute is a core designed property of the system. Partial ($\ominus$) indicates it is present but limited in scope, incidental, or a by-product of another objective. Absent ($\circ$) indicates it is not provided. Ratings reflect the systems as surveyed and would require revision as products change; per advisor guidance, the criteria set and ratings are to be finalised once the project scope is fixed.

\begin{landscape}
\begin{scriptsize}
\setlength{\tabcolsep}{3pt}
\renewcommand{\arraystretch}{1.3}
\begin{longtable}{@{}p{3.1cm} *{7}{p{2.55cm}}@{}}
\caption{Comparison of system classes against the derived criteria.}\\
\toprule
\textbf{Attribute} & \textbf{Authored simulation (Forage)} & \textbf{Generated simulation (EntryLevel, Anthropos)} & \textbf{Interview practice (careersim.ai etc.)} & \textbf{Generative narrative games (Hidden Door etc.)} & \textbf{Assessment \& orientation (O\textasteriskcentered NET, VN market)} & \textbf{Research prototypes (CareerSim, CareerPooler)} & \textbf{Proposed system} \\
\midrule
\endhead
C1 \newline Task authenticity & {\small$\bullet$}\  employer-authored real work products & {\small$\bullet$}\  generated role tasks & {\small$\ominus$}\  interview talk, not work & {\small$\circ$}\  fictional worlds & {\small$\circ$}\  inventory only & {\small$\ominus$}\  trajectory, not tasks & {\small$\bullet$}\  tasks drawn from occupational taxonomies \\
C2 \newline Candour of occupational representation & {\small$\bullet$}\  includes unglamorous tasks & {\small$\ominus$}\  varies by generation & {\small$\circ$}\  not its purpose & {\small$\circ$}\  no real occupations & {\small$\ominus$}\  text lists difficulties & {\small$\ominus$}\  events skew narrative & {\small$\bullet$}\  candour is a design requirement \\
C3 \newline Adaptivity to the learner & {\small$\circ$}\  fixed content, identical on repeat & {\small$\bullet$}\  parameterised generation & {\small$\bullet$}\  adaptive personas & {\small$\bullet$}\  fully generative & {\small$\ominus$}\  branching by score & {\small$\bullet$}\  generated per user & {\small$\bullet$}\  seeded generation per user \\
C4 \newline Validity of competence measurement & {\small$\circ$}\  model answer, no judgment & {\small$\ominus$}\  behavioural evidence, unvalidated & {\small$\ominus$}\  scored, conversational only & {\small$\circ$}\  none & {\small$\circ$}\  self-report only & {\small$\circ$}\  self-report outcomes & {\small$\ominus$}\  designed; validity to be established \\
C5 \newline Quality of feedback and guidance & {\small$\ominus$}\  exemplar, not response-specific & {\small$\ominus$}\  automated tutor & {\small$\bullet$}\  per-dimension coaching & {\small$\circ$}\  none & {\small$\ominus$}\  generic occupation reports & {\small$\ominus$}\  reflective prompts & {\small$\bullet$}\  result, feedback and guidance separated \\
C6 \newline Sustained engagement & {\small$\circ$}\  one-pass certificate & {\small$\circ$}\  discrete projects & {\small$\circ$}\  practice utility & {\small$\bullet$}\  high replayability & {\small$\circ$}\  single sitting & {\small$\ominus$}\  short study exposure & {\small$\ominus$}\  designed; retention untested \\
C7 \newline Support for pathway navigation & {\small$\circ$}\  isolated simulations & {\small$\circ$}\  unconnected projects & {\small$\circ$}\  single-role focus & {\small$\ominus$}\  world, not career, traversal & {\small$\bullet$}\  ranked occupational lists & {\small$\ominus$}\  linear life stages & {\small$\bullet$}\  competence-linked adjacency \\
C8 \newline Scalability and currency of content & {\small$\circ$}\  commissioned per employer & {\small$\bullet$}\  generated on demand & {\small$\bullet$}\  generated on demand & {\small$\bullet$}\  generated on demand & {\small$\ominus$}\  periodic database updates & {\small$\ominus$}\  prototype scale & {\small$\bullet$}\  generation with expert validation \\
\bottomrule
\end{longtable}
\end{scriptsize}

\noindent\scriptsize{$\bullet$~Full --- a core designed property \quad $\ominus$~Partial --- limited, incidental, or a by-product \quad $\circ$~Absent. Annotations state the basis for each rating. Ratings for the proposed system denote design intent; those marked $\ominus$ are targeted but not yet empirically established.}
\end{landscape}

\subsubsection{Gap and positioning}

The assessment indicates no systematic deficiency in any single system so much as a consistent failure to combine attributes that have been demonstrated separately. Authored simulation platforms achieve task authenticity and candour but cannot adapt or scale, and do not assess. Generated simulation systems achieve authenticity, adaptivity and scale, and where they assess behaviour they do so for selection rather than exploration, without pathway structure. Generative narrative systems achieve adaptivity and sustained engagement at high quality while making no representational commitment to real occupations. Assessment and orientation instruments --- including the entire Vietnamese market surveyed --- achieve pathway navigation and scale while omitting performance altogether. Research prototypes achieve adaptivity and engagement but stop short of task performance and competence measurement, and their evaluations do not support accuracy claims.

The consequent gap is not a missing feature but an unrealised combination: performed, candid occupational tasks that adapt to the individual, from which competence is measured with validated automated judgment, connected by a structure that relates occupations through the competencies they share, and maintained through expert authoring. Each constituent has been demonstrated --- task authenticity by Forage, generation and behavioural evidence by EntryLevel and Anthropos, coherent generative narrative by Hidden Door, pathway structure by occupational taxonomies, engagement by the research prototypes --- and their integration has not.

The present project addresses this combination. Its design adopts specific solutions identified above: bounded authored components as the means of maintaining generative consistency, discrete self-contained task structure for scenario construction, multi-role scenario composition for eliciting interpersonal competencies, balanced event generation for representing career uncertainty, and prospective framing at the entry point. Its contribution beyond assembly lies in the assessment and evaluation layer, where the requirements established in Section 2.5 --- chance-corrected agreement, cross-family evaluation, difficulty stratification --- are applied to a class of system in which such validation has not previously been reported. The claims the project may make are correspondingly bounded by Section 2.8: engagement and clarity effects are supported by existing evidence, whereas competence gain is a hypothesis to be tested against the reference set rather than a property the design confers.
